\documentclass[11pt,a4paper]{article}

% ── Packages (order matters: hyperref LAST) ─────────────────────────────────
\usepackage[utf8]{inputenc}
\usepackage[T1]{fontenc}
\usepackage{lmodern}
\usepackage{microtype}
\usepackage{amsmath,amssymb,amsthm}
\usepackage{mathtools}
\usepackage{graphicx}
\usepackage{booktabs}
\usepackage{longtable}
\usepackage{array}
\usepackage[margin=1in]{geometry}
\usepackage{algorithm}
\usepackage{algpseudocode}
\usepackage{subcaption}
\usepackage{xcolor}
\usepackage{natbib}
\usepackage{multirow}
\usepackage[hidelinks]{hyperref}   % load BEFORE cleveref (cleveref patches it)
\usepackage{cleveref}
\usepackage{doi}                    % render DOI fields as clickable links (after hyperref)

% ── Theorem-like environments ───────────────────────────────────────────────
\theoremstyle{plain}
\newtheorem{theorem}{Theorem}
\newtheorem{proposition}{Proposition}
\newtheorem{lemma}{Lemma}
\theoremstyle{remark}
\newtheorem{remark}{Remark}
\crefname{proposition}{Proposition}{Propositions}
\crefname{theorem}{Theorem}{Theorems}
\crefname{lemma}{Lemma}{Lemmas}
\crefname{remark}{Remark}{Remarks}

% ── Custom macros (project-specific — edit freely) ──────────────────────────
% \newcommand{\R}{\mathbb{R}}

% ── Figure placement discipline ─────────────────────────────────────────────
\usepackage{float}

\title{TODO: Paper Title\\Second Line if Needed}
\author{First Author\thanks{Affiliation, \texttt{email@example.com}}
  \and Second Author\thanks{Affiliation}}
\date{\today}

\begin{document}
\maketitle

\begin{abstract}
TODO: 150–250 word abstract. State the problem, the contribution, the
evidence, and the significance. Avoid citations here.
\end{abstract}

\tableofcontents
\newpage

%=============================================================================
\section{Introduction}
\label{sec:intro}
%=============================================================================

TODO: Motivation, gap, contributions (as a numbered list), roadmap.

Our contributions are:
\begin{enumerate}
  \item TODO
  \item TODO
  \item TODO
\end{enumerate}

The remainder of this paper is organized as follows.
\Cref{sec:background} reviews related work.
\Cref{sec:methods} presents the approach.
\Cref{sec:results} reports results.
\Cref{sec:discussion} discusses threats to validity.
\Cref{sec:conclusion} concludes.

%=============================================================================
\section{Background and Related Work}
\label{sec:background}
%=============================================================================

TODO. Example citation~\citep{author2024key}.

%=============================================================================
\section{Methods}
\label{sec:methods}
%=============================================================================

TODO.

% Example figure — copy this block for every figure.
\begin{figure}[t]
  \centering
  \includegraphics[width=0.85\textwidth]{figures/fig1_example.pdf}
  \caption{TODO: caption. Explain what the reader is looking at and what
    they should take away. Reference the generator script:
    \texttt{report/generate\_figures.py::plot\_example}.}
  \label{fig:example}
\end{figure}

%=============================================================================
\section{Results}
\label{sec:results}
%=============================================================================

TODO.

%=============================================================================
\section{Discussion}
\label{sec:discussion}
%=============================================================================

TODO. Threats to validity, limitations, generalization.

%=============================================================================
\section{Conclusion}
\label{sec:conclusion}
%=============================================================================

TODO.

%=============================================================================
% Bibliography
%=============================================================================
\bibliographystyle{plainnat}
\bibliography{references}

\end{document}
